\clearpage\thispagestyle{empty}
\pdfbookmark[0]{Quelle zu Juno Elemente VIII}{juno1808viii-sources}
\begin{center}\textbf{Quelle zu Juno Elemente VIII}\end{center}
\noindent Diese KI-Übertragung ins heutige Deutsch folgt der Mitteilung vom 25. August 1808 in Gauß’ \href{https://gdz.sub.uni-goettingen.de/id/PPN236007467}{\textit{Werke}, Bd.~VI} (Göttingen, 1874), S.~306--307. Der Erstdruck erschien unter der Überschrift \textit{Göttingen} in den \textit{Göttingischen gelehrten Anzeigen}, Stück 136, S.~\href{https://manifests.sub.uni-goettingen.de/iiif/image/gdz:PPN319721507_1808:00001359/full/1600,/0/default.jpg}{1353}, \href{https://manifests.sub.uni-goettingen.de/iiif/image/gdz:PPN319721507_1808:00001360/full/1600,/0/default.jpg}{1354} und \href{https://manifests.sub.uni-goettingen.de/iiif/image/gdz:PPN319721507_1808:00001361/full/1600,/0/default.jpg}{1355}.
\par\medskip
Die Mitteilung in der Werkausgabe verweist auf die Beobachtungstabelle im Brief vom 6. August. Der Erstdruck der Mitteilung gibt für die Rektaszension der Juno am 30. Juli die letzte Ziffer 0. Der Erstdruck des Briefs vom 6. August und die Werkausgabe geben dort 9. Die Mitteilung und der spätere Brief geben auch verschiedene Werte für die geozentrische Breite bei der Opposition. Die beiden Lesarten der Werkausgabe bleiben erhalten. KI-Übertragung und KI-Satz.
\clearpage
